% introduction_final.tex — Introduction for "The Thinking Tax"
% Loaded by main_final.tex via % introduction_final.tex — Introduction for "The Thinking Tax"
% Loaded by main_final.tex via % introduction_final.tex — Introduction for "The Thinking Tax"
% Loaded by main_final.tex via % introduction_final.tex — Introduction for "The Thinking Tax"
% Loaded by main_final.tex via \input{sections/introduction_final}

\begin{figure}[t]
    \centering
    \includegraphics[width=\linewidth]{fig_nothink_vs_thinking.pdf}
    \caption{\textbf{The Thinking Tax.}
    Non-thinking mode (blue) dramatically outperforms thinking mode (orange) at every matched token budget on GSM8K (Qwen3-8B).
    At the critical operating point, \texttt{nothink@256} achieves \textbf{87.5\%} accuracy with only 146 average tokens, while \texttt{think@512} reaches just \textbf{65.2\%} with 460 tokens---a \textbf{+22.3\,pp} accuracy advantage at \textbf{3.2$\times$} fewer tokens.
    The gap widens further for the 27B model (Table~\ref{tab:model-size-scaling}), whose longer reasoning chains are more severely truncated.}
    \label{fig:nothink-vs-thinking}
\end{figure}

Chain-of-thought (CoT) reasoning~\citep{wei2022chain} has become the dominant paradigm for improving LLM performance on complex tasks.
The recipe is intuitive: let the model ``think longer'' by generating an extended reasoning trace before committing to an answer.
This paradigm, operationalized in systems like OpenAI o1~\citep{openai2024o1}, QwQ~\citep{qwq2024}, and DeepSeek-R1~\citep{deepseekr1}, has produced impressive gains on mathematical reasoning benchmarks and is increasingly viewed as a reliable lever for scaling test-time performance~\citep{snell2024scaling}.
The implicit promise is simple: more thinking tokens, more accuracy.

We show that this promise breaks down in a surprisingly broad range of practical settings.

\paragraph{The thinking tax.}
Through systematic experiments on GSM8K with Qwen3-8B and Qwen3.5-27B, we uncover a striking phenomenon that we term the \emph{thinking tax}: \textbf{at every tested token budget, non-thinking mode outperforms thinking mode, often by a wide margin} (Figure~\ref{fig:nothink-vs-thinking}).
The headline result is stark: non-thinking mode with a 256-token budget achieves \textbf{87.5\%} accuracy using only 146 average tokens, while thinking mode with a \emph{twice-larger} 512-token budget reaches just \textbf{65.2\%}---requiring 460 average tokens.
This is not a marginal difference: non-thinking is \textbf{22.3 percentage points more accurate} while consuming \textbf{3.2$\times$ fewer tokens}.

The root cause is \emph{truncation waste}: thinking mode generates verbose reasoning chains that are truncated before reaching the answer.
While the mechanism may seem obvious in hindsight, neither its \emph{magnitude} (a 69.5\,pp gap at budget 256 on the full test set) nor its \emph{inverse scaling with model size} are predicted by simple format-overhead arguments.
The result is the worst of both worlds: the model spends its entire token budget on intermediate reasoning but never produces an answer, yielding 0\% accuracy on truncated samples by construction.
Non-thinking mode, by contrast, proceeds directly to the answer, using tokens far more efficiently and naturally stopping well within moderate budgets.

\paragraph{The thinking tax scales with model size.}
The problem is not confined to small models---it gets \emph{worse} as models grow larger.
The 27B model achieves only \textbf{18.3\%} accuracy in thinking mode at budget 512, compared to 8B's 65.2\%---a \textbf{46.9 percentage-point} collapse.
The mechanism is straightforward: larger models generate longer, more elaborate reasoning chains.
At budget 512, only \textbf{0.7\%} of 27B thinking traces produce a complete answer (vs.\ 6.1\% for 8B), meaning 99.3\% of responses are truncated projections that never reach a final answer.
The thinking tax at budget 512 is \textbf{2.6$\times$ larger} for 27B than for 8B (Table~\ref{tab:model-size-scaling}).
This inverse scaling property has troubling implications: as frontier models grow and their reasoning chains become longer and more complex, the practical cost of thinking mode under realistic deployment budgets will only increase.

\paragraph{Natural stopping as a confidence oracle.}
Our analysis reveals an intriguing structure within thinking mode's failures.
Among the minority of samples where the model's reasoning chain naturally terminates within the budget (i.e., the model ``decides it is done thinking''), accuracy is \textbf{96.3\%} for 8B at budget 512.
This natural stopping behavior functions as a remarkably accurate confidence oracle: when the model finishes thinking on its own, it is almost always correct.
The problem is that most samples \emph{never} reach this point within practical budgets.
This observation suggests a clean decomposition: some problems are genuinely easy for thinking mode (the model finishes quickly and is correct), while others are genuinely hard (the model cannot finish within budget).
The question is how to exploit this structure.

\paragraph{TOWN: Think Only When Needed.}
Building on the above insights, we propose \method{} (\textbf{\fullmethod{}}), a simple two-stage inference cascade:
\begin{enumerate}[nosep,leftmargin=*]
    \item \textbf{Stage 1: Non-thinking default.} Run the query in non-thinking mode with a moderate budget (e.g., 256 tokens). If the model produces a confident answer (which it does for ${\sim}89\%$ of GSM8K problems), accept it immediately.
    \item \textbf{Stage 2: Thinking fallback.} For the remaining hard problems where non-thinking mode fails or expresses low confidence, route to thinking mode with an extended budget, leveraging its strength on genuinely difficult instances.
\end{enumerate}
\method{} is motivated by two empirical observations: (i) non-thinking mode is both more accurate \emph{and} more efficient on the vast majority of problems; and (ii) thinking mode's natural-stop signal reliably identifies the minority of problems where extended reasoning genuinely helps.
On the full GSM8K test set ($n{=}1{,}319$), \method{} achieves \textbf{90.9\%} accuracy at only 199 average tokens---surpassing both \texttt{think@512} (+25.7\,pp) and \texttt{nothink@256} (+3.4\,pp).

\paragraph{Contributions.}
Our work makes three contributions:

\begin{enumerate}[nosep,leftmargin=*]
    \item \textbf{The Thinking Tax} (\S\ref{sec:thinking-tax}).
    We provide the first systematic empirical demonstration that non-thinking mode outperforms thinking mode at every tested token budget (128, 256, 512) on GSM8K, for both 8B and 27B models.
    At the most practically relevant operating point, non-thinking achieves +22.3\,pp higher accuracy at 3.2$\times$ fewer tokens.
    We trace the root cause to truncation waste and quantify the phenomenon through token utilization, natural-stop rates, and per-sample accuracy decomposition.

    \item \textbf{The Thinking Tax Scales with Model Size} (\S\ref{sec:finding:model-size}).
    We show that the thinking tax grows \emph{inversely} with model capability: the 27B model's thinking tax at budget 512 is 2.6$\times$ larger than the 8B model's, driven by longer reasoning chains and extremely low answer-completion rates (0.7\% vs.\ 6.1\%).
    This finding challenges the assumption that larger models will naturally overcome CoT's inefficiencies and has direct implications for the economics of deploying reasoning-augmented LLMs.

    \item \textbf{\method{}: Think Only When Needed} (\S\ref{sec:method}).
    We propose a training-free, two-stage cascade that defaults to non-thinking mode and selectively routes hard problems to thinking mode.
    \method{} achieves \textbf{90.9\%} accuracy at only 199 average tokens---\textbf{+25.7\,pp} over thinking@512 at 2.4$\times$ fewer tokens---on the full GSM8K test set, demonstrating that the question is not \emph{how long} to think, but \emph{whether} to think at all.
\end{enumerate}


\begin{figure}[t]
    \centering
    \includegraphics[width=\linewidth]{fig_nothink_vs_thinking.pdf}
    \caption{\textbf{The Thinking Tax.}
    Non-thinking mode (blue) dramatically outperforms thinking mode (orange) at every matched token budget on GSM8K (Qwen3-8B).
    At the critical operating point, \texttt{nothink@256} achieves \textbf{87.5\%} accuracy with only 146 average tokens, while \texttt{think@512} reaches just \textbf{65.2\%} with 460 tokens---a \textbf{+22.3\,pp} accuracy advantage at \textbf{3.2$\times$} fewer tokens.
    The gap widens further for the 27B model (Table~\ref{tab:model-size-scaling}), whose longer reasoning chains are more severely truncated.}
    \label{fig:nothink-vs-thinking}
\end{figure}

Chain-of-thought (CoT) reasoning~\citep{wei2022chain} has become the dominant paradigm for improving LLM performance on complex tasks.
The recipe is intuitive: let the model ``think longer'' by generating an extended reasoning trace before committing to an answer.
This paradigm, operationalized in systems like OpenAI o1~\citep{openai2024o1}, QwQ~\citep{qwq2024}, and DeepSeek-R1~\citep{deepseekr1}, has produced impressive gains on mathematical reasoning benchmarks and is increasingly viewed as a reliable lever for scaling test-time performance~\citep{snell2024scaling}.
The implicit promise is simple: more thinking tokens, more accuracy.

We show that this promise breaks down in a surprisingly broad range of practical settings.

\paragraph{The thinking tax.}
Through systematic experiments on GSM8K with Qwen3-8B and Qwen3.5-27B, we uncover a striking phenomenon that we term the \emph{thinking tax}: \textbf{at every tested token budget, non-thinking mode outperforms thinking mode, often by a wide margin} (Figure~\ref{fig:nothink-vs-thinking}).
The headline result is stark: non-thinking mode with a 256-token budget achieves \textbf{87.5\%} accuracy using only 146 average tokens, while thinking mode with a \emph{twice-larger} 512-token budget reaches just \textbf{65.2\%}---requiring 460 average tokens.
This is not a marginal difference: non-thinking is \textbf{22.3 percentage points more accurate} while consuming \textbf{3.2$\times$ fewer tokens}.

The root cause is \emph{truncation waste}: thinking mode generates verbose reasoning chains that are truncated before reaching the answer.
While the mechanism may seem obvious in hindsight, neither its \emph{magnitude} (a 69.5\,pp gap at budget 256 on the full test set) nor its \emph{inverse scaling with model size} are predicted by simple format-overhead arguments.
The result is the worst of both worlds: the model spends its entire token budget on intermediate reasoning but never produces an answer, yielding 0\% accuracy on truncated samples by construction.
Non-thinking mode, by contrast, proceeds directly to the answer, using tokens far more efficiently and naturally stopping well within moderate budgets.

\paragraph{The thinking tax scales with model size.}
The problem is not confined to small models---it gets \emph{worse} as models grow larger.
The 27B model achieves only \textbf{18.3\%} accuracy in thinking mode at budget 512, compared to 8B's 65.2\%---a \textbf{46.9 percentage-point} collapse.
The mechanism is straightforward: larger models generate longer, more elaborate reasoning chains.
At budget 512, only \textbf{0.7\%} of 27B thinking traces produce a complete answer (vs.\ 6.1\% for 8B), meaning 99.3\% of responses are truncated projections that never reach a final answer.
The thinking tax at budget 512 is \textbf{2.6$\times$ larger} for 27B than for 8B (Table~\ref{tab:model-size-scaling}).
This inverse scaling property has troubling implications: as frontier models grow and their reasoning chains become longer and more complex, the practical cost of thinking mode under realistic deployment budgets will only increase.

\paragraph{Natural stopping as a confidence oracle.}
Our analysis reveals an intriguing structure within thinking mode's failures.
Among the minority of samples where the model's reasoning chain naturally terminates within the budget (i.e., the model ``decides it is done thinking''), accuracy is \textbf{96.3\%} for 8B at budget 512.
This natural stopping behavior functions as a remarkably accurate confidence oracle: when the model finishes thinking on its own, it is almost always correct.
The problem is that most samples \emph{never} reach this point within practical budgets.
This observation suggests a clean decomposition: some problems are genuinely easy for thinking mode (the model finishes quickly and is correct), while others are genuinely hard (the model cannot finish within budget).
The question is how to exploit this structure.

\paragraph{TOWN: Think Only When Needed.}
Building on the above insights, we propose \method{} (\textbf{\fullmethod{}}), a simple two-stage inference cascade:
\begin{enumerate}[nosep,leftmargin=*]
    \item \textbf{Stage 1: Non-thinking default.} Run the query in non-thinking mode with a moderate budget (e.g., 256 tokens). If the model produces a confident answer (which it does for ${\sim}89\%$ of GSM8K problems), accept it immediately.
    \item \textbf{Stage 2: Thinking fallback.} For the remaining hard problems where non-thinking mode fails or expresses low confidence, route to thinking mode with an extended budget, leveraging its strength on genuinely difficult instances.
\end{enumerate}
\method{} is motivated by two empirical observations: (i) non-thinking mode is both more accurate \emph{and} more efficient on the vast majority of problems; and (ii) thinking mode's natural-stop signal reliably identifies the minority of problems where extended reasoning genuinely helps.
On the full GSM8K test set ($n{=}1{,}319$), \method{} achieves \textbf{90.9\%} accuracy at only 199 average tokens---surpassing both \texttt{think@512} (+25.7\,pp) and \texttt{nothink@256} (+3.4\,pp).

\paragraph{Contributions.}
Our work makes three contributions:

\begin{enumerate}[nosep,leftmargin=*]
    \item \textbf{The Thinking Tax} (\S\ref{sec:thinking-tax}).
    We provide the first systematic empirical demonstration that non-thinking mode outperforms thinking mode at every tested token budget (128, 256, 512) on GSM8K, for both 8B and 27B models.
    At the most practically relevant operating point, non-thinking achieves +22.3\,pp higher accuracy at 3.2$\times$ fewer tokens.
    We trace the root cause to truncation waste and quantify the phenomenon through token utilization, natural-stop rates, and per-sample accuracy decomposition.

    \item \textbf{The Thinking Tax Scales with Model Size} (\S\ref{sec:finding:model-size}).
    We show that the thinking tax grows \emph{inversely} with model capability: the 27B model's thinking tax at budget 512 is 2.6$\times$ larger than the 8B model's, driven by longer reasoning chains and extremely low answer-completion rates (0.7\% vs.\ 6.1\%).
    This finding challenges the assumption that larger models will naturally overcome CoT's inefficiencies and has direct implications for the economics of deploying reasoning-augmented LLMs.

    \item \textbf{\method{}: Think Only When Needed} (\S\ref{sec:method}).
    We propose a training-free, two-stage cascade that defaults to non-thinking mode and selectively routes hard problems to thinking mode.
    \method{} achieves \textbf{90.9\%} accuracy at only 199 average tokens---\textbf{+25.7\,pp} over thinking@512 at 2.4$\times$ fewer tokens---on the full GSM8K test set, demonstrating that the question is not \emph{how long} to think, but \emph{whether} to think at all.
\end{enumerate}


\begin{figure}[t]
    \centering
    \includegraphics[width=\linewidth]{fig_nothink_vs_thinking.pdf}
    \caption{\textbf{The Thinking Tax.}
    Non-thinking mode (blue) dramatically outperforms thinking mode (orange) at every matched token budget on GSM8K (Qwen3-8B).
    At the critical operating point, \texttt{nothink@256} achieves \textbf{87.5\%} accuracy with only 146 average tokens, while \texttt{think@512} reaches just \textbf{65.2\%} with 460 tokens---a \textbf{+22.3\,pp} accuracy advantage at \textbf{3.2$\times$} fewer tokens.
    The gap widens further for the 27B model (Table~\ref{tab:model-size-scaling}), whose longer reasoning chains are more severely truncated.}
    \label{fig:nothink-vs-thinking}
\end{figure}

Chain-of-thought (CoT) reasoning~\citep{wei2022chain} has become the dominant paradigm for improving LLM performance on complex tasks.
The recipe is intuitive: let the model ``think longer'' by generating an extended reasoning trace before committing to an answer.
This paradigm, operationalized in systems like OpenAI o1~\citep{openai2024o1}, QwQ~\citep{qwq2024}, and DeepSeek-R1~\citep{deepseekr1}, has produced impressive gains on mathematical reasoning benchmarks and is increasingly viewed as a reliable lever for scaling test-time performance~\citep{snell2024scaling}.
The implicit promise is simple: more thinking tokens, more accuracy.

We show that this promise breaks down in a surprisingly broad range of practical settings.

\paragraph{The thinking tax.}
Through systematic experiments on GSM8K with Qwen3-8B and Qwen3.5-27B, we uncover a striking phenomenon that we term the \emph{thinking tax}: \textbf{at every tested token budget, non-thinking mode outperforms thinking mode, often by a wide margin} (Figure~\ref{fig:nothink-vs-thinking}).
The headline result is stark: non-thinking mode with a 256-token budget achieves \textbf{87.5\%} accuracy using only 146 average tokens, while thinking mode with a \emph{twice-larger} 512-token budget reaches just \textbf{65.2\%}---requiring 460 average tokens.
This is not a marginal difference: non-thinking is \textbf{22.3 percentage points more accurate} while consuming \textbf{3.2$\times$ fewer tokens}.

The root cause is \emph{truncation waste}: thinking mode generates verbose reasoning chains that are truncated before reaching the answer.
While the mechanism may seem obvious in hindsight, neither its \emph{magnitude} (a 69.5\,pp gap at budget 256 on the full test set) nor its \emph{inverse scaling with model size} are predicted by simple format-overhead arguments.
The result is the worst of both worlds: the model spends its entire token budget on intermediate reasoning but never produces an answer, yielding 0\% accuracy on truncated samples by construction.
Non-thinking mode, by contrast, proceeds directly to the answer, using tokens far more efficiently and naturally stopping well within moderate budgets.

\paragraph{The thinking tax scales with model size.}
The problem is not confined to small models---it gets \emph{worse} as models grow larger.
The 27B model achieves only \textbf{18.3\%} accuracy in thinking mode at budget 512, compared to 8B's 65.2\%---a \textbf{46.9 percentage-point} collapse.
The mechanism is straightforward: larger models generate longer, more elaborate reasoning chains.
At budget 512, only \textbf{0.7\%} of 27B thinking traces produce a complete answer (vs.\ 6.1\% for 8B), meaning 99.3\% of responses are truncated projections that never reach a final answer.
The thinking tax at budget 512 is \textbf{2.6$\times$ larger} for 27B than for 8B (Table~\ref{tab:model-size-scaling}).
This inverse scaling property has troubling implications: as frontier models grow and their reasoning chains become longer and more complex, the practical cost of thinking mode under realistic deployment budgets will only increase.

\paragraph{Natural stopping as a confidence oracle.}
Our analysis reveals an intriguing structure within thinking mode's failures.
Among the minority of samples where the model's reasoning chain naturally terminates within the budget (i.e., the model ``decides it is done thinking''), accuracy is \textbf{96.3\%} for 8B at budget 512.
This natural stopping behavior functions as a remarkably accurate confidence oracle: when the model finishes thinking on its own, it is almost always correct.
The problem is that most samples \emph{never} reach this point within practical budgets.
This observation suggests a clean decomposition: some problems are genuinely easy for thinking mode (the model finishes quickly and is correct), while others are genuinely hard (the model cannot finish within budget).
The question is how to exploit this structure.

\paragraph{TOWN: Think Only When Needed.}
Building on the above insights, we propose \method{} (\textbf{\fullmethod{}}), a simple two-stage inference cascade:
\begin{enumerate}[nosep,leftmargin=*]
    \item \textbf{Stage 1: Non-thinking default.} Run the query in non-thinking mode with a moderate budget (e.g., 256 tokens). If the model produces a confident answer (which it does for ${\sim}89\%$ of GSM8K problems), accept it immediately.
    \item \textbf{Stage 2: Thinking fallback.} For the remaining hard problems where non-thinking mode fails or expresses low confidence, route to thinking mode with an extended budget, leveraging its strength on genuinely difficult instances.
\end{enumerate}
\method{} is motivated by two empirical observations: (i) non-thinking mode is both more accurate \emph{and} more efficient on the vast majority of problems; and (ii) thinking mode's natural-stop signal reliably identifies the minority of problems where extended reasoning genuinely helps.
On the full GSM8K test set ($n{=}1{,}319$), \method{} achieves \textbf{90.9\%} accuracy at only 199 average tokens---surpassing both \texttt{think@512} (+25.7\,pp) and \texttt{nothink@256} (+3.4\,pp).

\paragraph{Contributions.}
Our work makes three contributions:

\begin{enumerate}[nosep,leftmargin=*]
    \item \textbf{The Thinking Tax} (\S\ref{sec:thinking-tax}).
    We provide the first systematic empirical demonstration that non-thinking mode outperforms thinking mode at every tested token budget (128, 256, 512) on GSM8K, for both 8B and 27B models.
    At the most practically relevant operating point, non-thinking achieves +22.3\,pp higher accuracy at 3.2$\times$ fewer tokens.
    We trace the root cause to truncation waste and quantify the phenomenon through token utilization, natural-stop rates, and per-sample accuracy decomposition.

    \item \textbf{The Thinking Tax Scales with Model Size} (\S\ref{sec:finding:model-size}).
    We show that the thinking tax grows \emph{inversely} with model capability: the 27B model's thinking tax at budget 512 is 2.6$\times$ larger than the 8B model's, driven by longer reasoning chains and extremely low answer-completion rates (0.7\% vs.\ 6.1\%).
    This finding challenges the assumption that larger models will naturally overcome CoT's inefficiencies and has direct implications for the economics of deploying reasoning-augmented LLMs.

    \item \textbf{\method{}: Think Only When Needed} (\S\ref{sec:method}).
    We propose a training-free, two-stage cascade that defaults to non-thinking mode and selectively routes hard problems to thinking mode.
    \method{} achieves \textbf{90.9\%} accuracy at only 199 average tokens---\textbf{+25.7\,pp} over thinking@512 at 2.4$\times$ fewer tokens---on the full GSM8K test set, demonstrating that the question is not \emph{how long} to think, but \emph{whether} to think at all.
\end{enumerate}


\begin{figure}[t]
    \centering
    \includegraphics[width=\linewidth]{fig_nothink_vs_thinking.pdf}
    \caption{\textbf{The Thinking Tax.}
    Non-thinking mode (blue) dramatically outperforms thinking mode (orange) at every matched token budget on GSM8K (Qwen3-8B).
    At the critical operating point, \texttt{nothink@256} achieves \textbf{87.5\%} accuracy with only 146 average tokens, while \texttt{think@512} reaches just \textbf{65.2\%} with 460 tokens---a \textbf{+22.3\,pp} accuracy advantage at \textbf{3.2$\times$} fewer tokens.
    The gap widens further for the 27B model (Table~\ref{tab:model-size-scaling}), whose longer reasoning chains are more severely truncated.}
    \label{fig:nothink-vs-thinking}
\end{figure}

Chain-of-thought (CoT) reasoning~\citep{wei2022chain} has become the dominant paradigm for improving LLM performance on complex tasks.
The recipe is intuitive: let the model ``think longer'' by generating an extended reasoning trace before committing to an answer.
This paradigm, operationalized in systems like OpenAI o1~\citep{openai2024o1}, QwQ~\citep{qwq2024}, and DeepSeek-R1~\citep{deepseekr1}, has produced impressive gains on mathematical reasoning benchmarks and is increasingly viewed as a reliable lever for scaling test-time performance~\citep{snell2024scaling}.
The implicit promise is simple: more thinking tokens, more accuracy.

We show that this promise breaks down in a surprisingly broad range of practical settings.

\paragraph{The thinking tax.}
Through systematic experiments on GSM8K with Qwen3-8B and Qwen3.5-27B, we uncover a striking phenomenon that we term the \emph{thinking tax}: \textbf{at every tested token budget, non-thinking mode outperforms thinking mode, often by a wide margin} (Figure~\ref{fig:nothink-vs-thinking}).
The headline result is stark: non-thinking mode with a 256-token budget achieves \textbf{87.5\%} accuracy using only 146 average tokens, while thinking mode with a \emph{twice-larger} 512-token budget reaches just \textbf{65.2\%}---requiring 460 average tokens.
This is not a marginal difference: non-thinking is \textbf{22.3 percentage points more accurate} while consuming \textbf{3.2$\times$ fewer tokens}.

The root cause is \emph{truncation waste}: thinking mode generates verbose reasoning chains that are truncated before reaching the answer.
While the mechanism may seem obvious in hindsight, neither its \emph{magnitude} (a 69.5\,pp gap at budget 256 on the full test set) nor its \emph{inverse scaling with model size} are predicted by simple format-overhead arguments.
The result is the worst of both worlds: the model spends its entire token budget on intermediate reasoning but never produces an answer, yielding 0\% accuracy on truncated samples by construction.
Non-thinking mode, by contrast, proceeds directly to the answer, using tokens far more efficiently and naturally stopping well within moderate budgets.

\paragraph{The thinking tax scales with model size.}
The problem is not confined to small models---it gets \emph{worse} as models grow larger.
The 27B model achieves only \textbf{18.3\%} accuracy in thinking mode at budget 512, compared to 8B's 65.2\%---a \textbf{46.9 percentage-point} collapse.
The mechanism is straightforward: larger models generate longer, more elaborate reasoning chains.
At budget 512, only \textbf{0.7\%} of 27B thinking traces produce a complete answer (vs.\ 6.1\% for 8B), meaning 99.3\% of responses are truncated projections that never reach a final answer.
The thinking tax at budget 512 is \textbf{2.6$\times$ larger} for 27B than for 8B (Table~\ref{tab:model-size-scaling}).
This inverse scaling property has troubling implications: as frontier models grow and their reasoning chains become longer and more complex, the practical cost of thinking mode under realistic deployment budgets will only increase.

\paragraph{Natural stopping as a confidence oracle.}
Our analysis reveals an intriguing structure within thinking mode's failures.
Among the minority of samples where the model's reasoning chain naturally terminates within the budget (i.e., the model ``decides it is done thinking''), accuracy is \textbf{96.3\%} for 8B at budget 512.
This natural stopping behavior functions as a remarkably accurate confidence oracle: when the model finishes thinking on its own, it is almost always correct.
The problem is that most samples \emph{never} reach this point within practical budgets.
This observation suggests a clean decomposition: some problems are genuinely easy for thinking mode (the model finishes quickly and is correct), while others are genuinely hard (the model cannot finish within budget).
The question is how to exploit this structure.

\paragraph{TOWN: Think Only When Needed.}
Building on the above insights, we propose \method{} (\textbf{\fullmethod{}}), a simple two-stage inference cascade:
\begin{enumerate}[nosep,leftmargin=*]
    \item \textbf{Stage 1: Non-thinking default.} Run the query in non-thinking mode with a moderate budget (e.g., 256 tokens). If the model produces a confident answer (which it does for ${\sim}89\%$ of GSM8K problems), accept it immediately.
    \item \textbf{Stage 2: Thinking fallback.} For the remaining hard problems where non-thinking mode fails or expresses low confidence, route to thinking mode with an extended budget, leveraging its strength on genuinely difficult instances.
\end{enumerate}
\method{} is motivated by two empirical observations: (i) non-thinking mode is both more accurate \emph{and} more efficient on the vast majority of problems; and (ii) thinking mode's natural-stop signal reliably identifies the minority of problems where extended reasoning genuinely helps.
On the full GSM8K test set ($n{=}1{,}319$), \method{} achieves \textbf{90.9\%} accuracy at only 199 average tokens---surpassing both \texttt{think@512} (+25.7\,pp) and \texttt{nothink@256} (+3.4\,pp).

\paragraph{Contributions.}
Our work makes three contributions:

\begin{enumerate}[nosep,leftmargin=*]
    \item \textbf{The Thinking Tax} (\S\ref{sec:thinking-tax}).
    We provide the first systematic empirical demonstration that non-thinking mode outperforms thinking mode at every tested token budget (128, 256, 512) on GSM8K, for both 8B and 27B models.
    At the most practically relevant operating point, non-thinking achieves +22.3\,pp higher accuracy at 3.2$\times$ fewer tokens.
    We trace the root cause to truncation waste and quantify the phenomenon through token utilization, natural-stop rates, and per-sample accuracy decomposition.

    \item \textbf{The Thinking Tax Scales with Model Size} (\S\ref{sec:finding:model-size}).
    We show that the thinking tax grows \emph{inversely} with model capability: the 27B model's thinking tax at budget 512 is 2.6$\times$ larger than the 8B model's, driven by longer reasoning chains and extremely low answer-completion rates (0.7\% vs.\ 6.1\%).
    This finding challenges the assumption that larger models will naturally overcome CoT's inefficiencies and has direct implications for the economics of deploying reasoning-augmented LLMs.

    \item \textbf{\method{}: Think Only When Needed} (\S\ref{sec:method}).
    We propose a training-free, two-stage cascade that defaults to non-thinking mode and selectively routes hard problems to thinking mode.
    \method{} achieves \textbf{90.9\%} accuracy at only 199 average tokens---\textbf{+25.7\,pp} over thinking@512 at 2.4$\times$ fewer tokens---on the full GSM8K test set, demonstrating that the question is not \emph{how long} to think, but \emph{whether} to think at all.
\end{enumerate}
